\subsection{The direction selected by the dodecaphase record}
\label{exc:k-direccion}

The record \(K\) has a geometric function complementing its
rational evaluation and support reading. Once the representation
of the twelve positions has been constructed, its components determine a
direction in the space of zero-mean variations. The reduction
\(12\to11\) eliminates the uniform mode; the following reduction
\(11\to10\) eliminates the direction selected by the record itself.
The two operations therefore have different antecedents.

To declare the chart used completely, consider the group
\(A_5\) of even permutations of five objects and the subgroup
\(C_5=\langle(1\,2\,3\,4\,5)\rangle\). The twelve left cosets
\(r_pC_5\) are ordered through these representatives, written as
image lists:
\[
\begin{array}{c|c@{\qquad}c|c}
p&r_p&p&r_p\\ \hline
1&(4,3,5,2,1)&7&(5,4,3,2,1)\\
2&(4,2,1,3,5)&8&(1,3,4,2,5)\\
3&(1,2,4,5,3)&9&(4,2,3,5,1)\\
4&(2,5,3,4,1)&10&(3,2,5,4,1)\\
5&(4,3,1,5,2)&11&(4,5,2,3,1)\\
6&(5,1,2,3,4)&12&(5,3,2,4,1).
\end{array}
\]
The action by left multiplication defines an orthogonal
representation \(\rho:A_5\to O(\RR^{12})\):
\(\rho(a)e_p=e_q\) when \(ar_pC_5=r_qC_5\).
This representation chart acts on the record already obtained;
group characters are not used to choose its twelve values.

Let \(\chi_3\) be the values of the three-dimensional character
\[
\begin{array}{c|ccccc}
\text{class}&1&2^2&3&5_a&5_b\\ \hline
\chi_3&3&-1&0&(1+\sqrt5)/2&(1-\sqrt5)/2.
\end{array}
\]
The class \(5_a\) contains the displayed generator of \(C_5\), and
\(5_b\) contains its square. This convention distinguishes the two
conjugate three-dimensional sectors of the representation. Define
\begin{equation}
 P_3=\frac3{60}\sum_{a\in A_5}\chi_3(a^{-1})\rho(a),
 \qquad P_{11}=I_{12}-\frac1{12}J_{12}.
 \label{eq:k-proyectores-incidencia}
\end{equation}
Both matrices are determined by finite sums in
\(\QQ(\sqrt5)\). The appearance of this quadratic field belongs
to the realization of the representation, subsequent to regional
generation of \(\varphi\).

\begin{lemma}[Projection of the record onto the three-dimensional sector]
\label{lem:k-sector-no-nulo}
One has
\[
 P_3=P_3^{\mathsf T}=P_3^2,\quad
 P_3\mathbf1=0,\quad \operatorname{tr}P_3=3,
 \qquad P_3P_{11}=P_{11}P_3=P_3.
\]
For the record \eqref{exc:k}, the vector
\begin{equation}
 u_K=P_3P_{11}K
 \label{eq:k-vector-direccional}
\end{equation}
is nonzero. Its exact norm is
\begin{equation}
 \lVert u_K\rVert^2
 =\frac{6638585+2275584\sqrt5}{20}>0.
 \label{eq:k-norma-direccion}
\end{equation}
\end{lemma}
\begin{proof}
The character sum in \eqref{eq:k-proyectores-incidencia} is the
isotypic projector. The real character and substitution \(a\mapsto a^{-1}\)
give symmetry. The multiplicity of the sector in the action on
\(A_5/C_5\) is the dimension of its vectors fixed by \(C_5\),
\[
 \frac15\left(3+2\frac{1+\sqrt5}{2}
                 +2\frac{1-\sqrt5}{2}\right)=1.
\]
Its rank is therefore three. Orthogonality to the trivial character
gives \(P_3\mathbf1=0\) and the two compositions with \(P_{11}\).
These identities may also be verified by multiplying the
matrices in the finite sum.

Finally, \(\sum K_p=6263\), and substituting the twelve
components into that same sum gives
\[
 \lVert u_K\rVert^2=K^{\mathsf T}P_3K
 =\frac1{20}\sum_{a\in A_5}
                \chi_3(a^{-1})K^{\mathsf T}\rho(a)K
 =\frac{6638585+2275584\sqrt5}{20}.
\]
This is an exact evaluation of sixty specified terms,
reproduced by the accompanying certificate using quadratic rational
arithmetic. The two positive coefficients prove nonvanishing.
\end{proof}

\begin{definition}[Dodecaphase direction and transversal complement]
The direction selected by \(K\) in the declared chart is
\(\ell_K=\RR u_K\). Its transversal complement within
\(V_{11}=\mathbf1^\perp\) is
\[
 V_{10}(K)=V_{11}\cap u_K^\perp.
\]
\end{definition}

\begin{proposition}[Orthogonal reduction determined by \(K\)]
\label{prop:k-reduccion-diez}
The operator
\begin{equation}
 P_{10}(K)=P_{11}
       -\frac{u_Ku_K^{\mathsf T}}{\lVert u_K\rVert^2}
 \label{eq:k-proyector-diez}
\end{equation}
is the rank-ten orthogonal projector onto \(V_{10}(K)\). One has
\[
 V_{11}=\ell_K\mathbin{\oplus^{\perp}}V_{10}(K),
 \qquad P_{10}P_{11}=P_{10},\quad P_{10}u_K=0.
\]
\end{proposition}
\begin{proof}
Let \(E_K=u_Ku_K^{\mathsf T}/\lVert u_K\rVert^2\). The
proved nonvanishing guarantees that it is defined. It is symmetric,
\(E_K^2=E_K\), has rank one, and, since \(u_K\in V_{11}\),
\(P_{11}E_K=E_KP_{11}=E_K\). Hence
\[
 (P_{11}-E_K)^2=P_{11}-2E_K+E_K=P_{11}-E_K.
\]
Its image is \(V_{11}\cap u_K^\perp\), and its trace is
\(11-1=10\). The decomposition and remaining identities are
obtained by projection onto the two orthogonal summands.
\end{proof}

This result specifies a function of \(K\) not exhausted by its designation as
a record: in addition to reversibly representing oriented
data and selecting the tetrad \(B_K\), it determines a
linear polarization. The direction depends on the complete record
through \(P_3K\), not only on the four points of its support.
The order record, representation, direction, and
dimensional reduction has thus been maintained.

The transformation has a precise covariance law. If the chart is
changed simultaneously by an orthogonal permutation matrix \(S\),
\(K'=SK\) and \(P'_3=SP_3S^{-1}\), then
\[
 u_{K'}'=Su_K,\qquad P'_{10}(K')=SP_{10}(K)S^{-1}.
\]
Moreover, replacing \(u_K\) by \(a u_K\), with \(a\ne0\),
does not change \(E_K\). The direction retains a polarization and does not
by itself allow recovery of either the amplitude of \(u_K\) or the
complete record. Reversibility of \(\kappa\leftrightarrow K\)
remains in its own domain.

The geometric realization established here is a reduction of inner-product
spaces. Its subsequent use as a compact direction
requires transporting this line to the corresponding physical realization.
Section \ref{exc:pantallas-dualidad} sets out the algebraic map
keeping this reduction and the lattice screen coordinated.
